
\section{Conclusion}\label{sec:conclusion}\label{sec:disc-summary}

This thesis set out to validate a taxonomy and ended by measuring a
judge. Restating the contributions of
Section~\ref{sec:research-goal} as what was demonstrated rather than
what was attempted:

\textbf{What was demonstrated.}
On a corpus with a verifiable answer key, an LLM judge's pairwise
verdicts track answer correctness rather than reasoning quality. That
is established by five measurements with mutually independent failure
modes: an agreement rate with the key of $88.8\%$ on discriminable
items; a tie rate that triples exactly where the key stops
discriminating; an inversion by trace presence that \emph{strengthens}
to $97.0\%$ against $83.2\%$ once capability is controlled; a
rubric-scoped against rubric-free contrast that leaves $99.4\%$ of
winners identical; and a paired partition against no-partition contrast
that is indistinguishable on its high-power statistic
($94.2\%$ against $94.7\%$).
A construction method was built and characterised as an instrument: it
reaches a certified structural fixed point at iteration 10, its
acceptance rule is calibrated in units of its own objective's paired
standard error against a $13.1\times$ spread in that error, and its
separation threshold is positioned between two measured nulls.
Rubrics induced on its cells were shown to be specific to those cells
against a size-matched randomised null at
$p = 1.21 \times 10^{-6}$, with all four pre-registered discriminators
firing in their predicted directions.
And the premise of this thesis's original second research question ---
that cells carry different true model rankings --- was tested and
refused under pre-registration, before the arena ran.

\textbf{Answering the research questions.}
RQ1 asked whether answer-key-blind, domain-conditioned pairwise
judgments recover per-domain capability rankings, and what decides the
verdicts. They recover the ground-truth ordering to within about one
adjacent transposition, and what decides the verdicts is answer
correctness; on the pilot roster, half of that recovery was purchased
by a response-format artifact rather than earned by judgment.
RQ2 asked whether adapted, cell-scoped grouping enables better judging
than a generic partition-free rubric on the same questions. On Math ---
designated in advance as the null arm, because the roster's models
order identically there and globally --- the two arms are
indistinguishable on answer-key agreement, and the rank-correlation gap
that appears under one tie convention collapses to the registered
boundary under the other. Under the registered rule that gap counts as
a \emph{failed} null in favour of the partition-free arm; the
tie-scoring attribution is post hoc, and the confirmatory design fixes
the convention in advance. The corresponding test on law, the domain the
offline screen selects, is outstanding.

\textbf{What the scope licenses, and what it does not.}
These are results about one multiple-choice corpus, one judge model,
one or two domains, rosters of 8 to 12 models, and a single seed. They
license the claim that on a corpus of this kind the judging mechanism
can bypass the partition, and they license the negative result about
geometry-adapted evaluation partitioning for a roster of this
capability spread. They do \emph{not} license a claim that cell-scoped
rubrics fail to improve judging in general: the regime in which a
rubric has been shown able to matter --- where correctness is not
checkable --- has not been tested here. Nor do they license any
absolute rank-fidelity figure quoted to three decimals
(Section~\ref{sec:res-tie-policy}); only paired contrasts within a run
are reported as such.

\textbf{What the reframing contributes.}
The evaluative value of a partition is a claim about a judge's
generalisation surface, and it decomposes into three separately
falsifiable links: coherent cells, specific rubrics, better judgments.
Stated that way, this project's negative results become locatable ---
link~1 holds, link~2 holds, and link~3 fails at its premise on this
corpus and at its mechanism on this judging paradigm. Stated the other
way, as a single claim that adapted partitioning improves ranking
fidelity, the same results read as a project that did not work. The
decomposition is not a post-hoc rescue; it was written down and dated
before the runs it now organises, and it is the difference between a
result and a shrug.

\section{Outlook and Future Directions}\label{sec:outlook-and-future-directions}\label{sec:future-work}\label{sec:disc-future}

In priority order, with the reason each sits where it does.

\begin{enumerate}

\item \textbf{A rule-identifier citation scheme in the verdict schema.}
Section~\ref{sec:res-generic-judge} shows that the judge reproduces a
cell rubric's vocabulary; it cannot show that the judge \emph{applies}
the rubric's criteria, because lexical overlap cannot distinguish the
two. Requiring the verdict to name the rule it relied on settles that
directly --- a cited rule identifier cannot arise from topical
coincidence --- and it counts, for free, how many verdicts cite no rule
at all and are therefore decided on priors. It is first on this list
because it changes the verdict schema and therefore must precede any
further arena run.

\item \textbf{The options-blind arm.} The same held-out questions, the
same cell rubrics, with the multiple-choice options withheld from the
judge. This is the demonstration that Section~\ref{sec:disc-rq1}
currently has only as an inference, and its predictions are already
fixed in writing per stratum. It is \emph{blocked} until the roster is
trace-complete: without options and without stored reasoning, an
untraced-against-untraced comparison presents the judge with a question
and two bare letters, which is uninterpretable rather than a
chance-level null. The corpus repair of
Section~\ref{sec:res-capture-gap} removes that block for most of the
corpus, and re-running generation with output capture removes it for
the rest.

\item \textbf{The law arena on the 11-model band} (Run~C). Law is the
one domain that reorders significantly under both roster policies at
full question coverage, and it is therefore the domain where the
partition has something to be right about. Without it,
Section~\ref{sec:res-law} reports the screen as a \emph{bound} on what
a partition could show rather than as a test of whether it shows it.

\item \textbf{A roster clustered within about ten accuracy points.} The
granularity null of Section~\ref{sec:res-granularity} was measured on a
roster spanning $54$ accuracy points, where global capability ordering
dominates by construction. A tightly clustered roster is the version of
that question that could still come out positive, and it is untested.

\item \textbf{Generalisation to all 14 domains at matched statistical
power}, rather than one or two selected domains at high power.

\item \textbf{A general trace-presence threshold in the corpus
validator.} The current exclusion policy is a curated list of ten model
names. A failure branch on the trace-presence check at, say, $0.90$
turns it into a threshold that generalises to models nobody has
inspected. It is a one-line change and it makes the limitation of
Section~\ref{subsec:disc-external-validity} actionable rather than
merely honest.

\item \textbf{A Davidson tie model.} Section~\ref{sec:res-tie-policy}
shows that the choice between half-weighting and dropping ties is not
neutral on this pipeline. Fitting tie propensity as a parameter
replaces a convention that currently changes conclusions with an
estimate that can be reported with an interval.

\item \textbf{A direct estimate of judge memorisation}, by verdict
consistency under paraphrased questions or by a probe set published
after the judge model's training cutoff. The answer-key-blind protocol
cannot defend against contamination that sits in the judge's weights.

\item \textbf{Exposing the routing component as a deployable service.}
This thesis builds the substrate a router would consume and does not
deploy one; the routing rule exists and is exercised read-only during
each arena run.

\end{enumerate}
